\documentclass[11pt]{article}

\usepackage{amsmath,amssymb}
\usepackage{xurl}
\usepackage[hidelinks]{hyperref}
\setlength{\emergencystretch}{2em}

% Shared commands for notes in docs/paper-gaps/.
%
% Two halves.  The link commands below are project identity: OWNER/REPO is the
% placeholder that scripts/bootstrap_project.py replaces with this project's
% GitHub slug and Pages site.  Everything after them is NOTATION, and notation
% belongs to the source paper: the definitions here are an example set, and a
% project replaces them with the macros of its own paper's style file.  The one
% rule that never changes: a command defined here must mean exactly what the
% source paper's macro means, or a note silently says something else.

\newcommand{\Lean}{\textsc{Lean}}
\newcommand{\leanid}[1]{\textsf{#1}}
\newcommand{\ghissue}[1]{\href{https://github.com/OWNER/REPO/issues/#1}{GitHub issue \##1}}
\newcommand{\ghpr}[1]{\href{https://github.com/OWNER/REPO/pull/#1}{PR \##1}}
% Local issue tracker (issues/NNNN-*.md in this repository); no remote URL.
\newcommand{\localissue}[1]{issue \##1 (\path{issues/})}
\newcommand{\doclink}[1]{\href{https://OWNER.github.io/REPO/docs/#1.html}{\path{#1}}}

% --- Notation: replace with the source paper's own macros --------------------
% \F, \N, \C, \R, \tr, \ot, \eps, \E, \bx, \bu, \bv, \by

\newcommand{\F}{\mathbb{F}}
\newcommand{\N}{\mathbb{N}}
\newcommand{\C}{\mathbb{C}}
\newcommand{\R}{\mathbb{R}}
\newcommand{\tr}{\mathrm{tr}}
\newcommand{\eps}{\epsilon}
\newcommand{\ot}{\otimes}
\newcommand{\E}{\mathop{\bf E\/}}

% bold random variables
\newcommand{\bu}{\boldsymbol{u}}
\newcommand{\bv}{\boldsymbol{v}}
\newcommand{\bx}{{\boldsymbol{x}}}
\newcommand{\by}{\boldsymbol{y}}

% T1 so literal < and > in \texttt placeholders typeset as themselves
\usepackage[T1]{fontenc}

% bra-ket
\usepackage{braket}

% --- Example structured spaces ---
\newcommand{\polyfunc}[3]{\mathcal{P}(#1, #2, #3)}
\newcommand{\polymeas}[3]{\mathrm{PolyMeas}(#1, #2, #3)}
\newcommand{\polysub}[3]{\mathrm{PolySub}(#1, #2, #3)}

% Approximation relations (paper uses \approx_\eps and \simeq_\zeta)
\newcommand{\approxeps}{\approx_{\varepsilon}}
\newcommand{\simgeq}{\simeq_{\zeta}}

\renewcommand{\ghissue}[1]{%
  \href{https://github.com/Dengnifer/MIPStarRE-A/issues/#1}{GitHub issue \##1}}

\title{The Cross-Basis Phase in the Extracted Pauli Observables}
\author{}
\date{2026-09-22}

\begin{document}
\maketitle

\section{At a glance}

\begin{itemize}
  \item \textbf{Difficulty.} Low.  The source drops a nontrivial phase when the
  two binary-basis indices are distinct.
  \item \textbf{Estimated weight.} One exact multiplication calculation for
  generalized Pauli observables; no robustness estimate changes.
  \item \textbf{Mathlib/project split.} Finite-sum bilinearity and the field
  trace are standard algebra.  The extracted-observable product form and the
  project's phase convention are project-local.
  \item \textbf{Key mathematical inputs.} The generalized Pauli commutation
  relation and nondegeneracy of the trace pairing.  The local correction is
  tracked in \ghissue{19}.
\end{itemize}

\section{Key theorem forms}

Let $e_0,\ldots,e_{t-1}$ be the fixed self-dual basis of $\mathbb F_q$ over
$\mathbb F_2$.  With the phase convention of the source, generalized Pauli
observables satisfy
\[
  \tau^X(a)\tau^Z(b)
  =(-1)^{\operatorname{tr}(a\mathbin{\cdot}b)}
   \tau^Z(b)\tau^X(a).
\]
Substituting $a=e_j\widetilde u$ and $b=e_{j'}\widetilde v$ gives the exact
relation
\[
  \widetilde X^j(\widetilde u)\widetilde Z^{j'}(\widetilde v)
  =(-1)^{\operatorname{tr}
      (e_je_{j'}(\widetilde u\mathbin{\cdot}\widetilde v))}
   \widetilde Z^{j'}(\widetilde v)\widetilde X^j(\widetilde u).
\]
This formula applies uniformly, whether or not $j=j'$.

\section{The false dichotomy in the source}

Lines~1438--1450 of
\path{references/qpbt-paper/14_analysis_of_the_pauli_basis_test.tex}
identify each extracted observable with a tensor product containing
$\tau^W(e_j\widetilde u)$.  Lines~1451--1456 then claim that the extracted
$X$- and $Z$-observables commute whenever $j\neq j'$
\cite[lines~1438--1456]{Ji2020MIPStar}.  Self-duality gives only
\[
  \operatorname{tr}(e_je_{j'})=0 \qquad (j\neq j').
\]
It does not imply
\[
  \operatorname{tr}(e_je_{j'}s)=0
  \qquad\text{for every }s\in\mathbb F_q.
\]

Indeed, for distinct $j,j'$, the element $c=e_je_{j'}$ is nonzero.  The trace
pairing is nondegenerate, so there is an $s\in\mathbb F_q$ for which
$\operatorname{tr}(cs)=1$.  Taking vectors whose first coordinates are $1$
and $s$, with all remaining coordinates zero, gives
$\widetilde u\mathbin{\cdot}\widetilde v=s$.  The two extracted observables
then anticommute.  This is a counterexample to the printed cross-basis
commutation assertion.

For a concrete admissible parameter tuple, take $(q,m,d)=(8,1,2)$.
The binary basis has three elements, so two distinct indices exist.  With
$c=e_je_{j'}$, take $s=c^{-1}$; then
$\operatorname{tr}(cs)=\operatorname{tr}(1)=1$, since the extension degree is
three.  There are $2^m=2$ register coordinates, enough for the vectors just
described.  The extracted observables are involutions, so their product is
invertible on the nonzero strategy space.  Anticommutation therefore cannot
also be commutation.  The counterexample is within the source domain and does
not depend on which projective polynomial-pair measurement was constructed.

When $\widetilde u\mathbin{\cdot}\widetilde v$ lies in the prime field, the
printed distinction is recovered from self-duality.  The chapter quantifies
over arbitrary vectors, however, so that special case does not repair its
statement.

The corrected formula follows directly from the product form: the two
sign-weighted sums of the same orthogonal projectors commute, and the remaining
Pauli factors contribute exactly the displayed trace phase.  This proof covers
equal indices, distinct indices, and zero register vectors.  No restriction of
the field-valued vectors or of their dot product is needed.

\section{Blueprint and formalization status}

Equation~\path{eq:tildew-twisted-commutation} and
Remark~\path{rem:qld-cross-phase} in
\path{blueprint/src/chapter/ch16_qpbt_extraction.tex} record the full phase and
note that the false dichotomy is not used later in the extraction argument.
The declaration \path{tildeObs_twisted_commutation} in
\path{MIPStarRE/QPBT/Extraction/Observables.lean} quantifies
$j,j'$ over the zero-based type
$\operatorname{Fin}(\texttt{P.model.basisDim})$ and concludes exactly the
uniform phase relation displayed above.  Its docstring carries the required
\textbf{Local fix} marker.  Thus the Lean statement agrees with the corrected
blueprint and deliberately differs from the false cross-basis commutation
claim in the paper; it introduces no extra hypothesis.

The declaration is now proved from the product form and generalized Pauli
commutation; this independent algebraic result is recorded in \ghissue{240}.
Issue~\#19 records the Stage~4.2 statement and this discrepancy. The
global polynomial-pair construction and the extraction theorems now have
separate completed proofs. They do not depend on the false cross-basis
dichotomy; the calculation here does not discharge the source-law obligations
recorded in \path{docs/paper-gaps/qpbt_ld-dimension-divisibility.tex}.

This correction changes neither the definition of the extracted observables
nor the source-faithful consistency and extraction conclusions.  It repairs an
auxiliary algebraic assertion that is independent of the robustness estimates.

\section{Retention and consumer comparison}

The printed assertion is retained as a proposition, not as a theorem.  Both
branches of the dichotomy are present: unconditional commutation for distinct
indices and the printed trace phase for equal indices.  The parameters,
projective strategy, and global measurement are the ambient objects already
fixed in the source; the vectors still range over all of $\mathbb F_q^{2^m}$.
Only the basis indexing changes from one-based to zero-based.\footnote{The
unasserted definition is
\leanid{MIPStarRE.QPBT.PrintedCrossBasisCommutationClaim} in
\path{MIPStarRE/QPBT/Extraction/PrintedClaims.lean}, tracked in
\ghissue{667}.  It is not a hypothesis of any theorem.  The corrected
blueprint node and its note citation are unchanged.}

The subsequent source argument has three relevant uses of the extracted
observables.  The consistency lemma uses their defining spectral sums and
same-basis measurement consistency.  The swap calculation at lines~1687--1713
uses the product form and commutation of the original generalized Pauli
operators, not the false dichotomy for extracted operators.  The EPR-state and
Pauli-measurement extraction, and then the soundness theorem, use those
consistency and conjugation conclusions.  Each use is preserved.  In the
blueprint, the dependence of the swap-conjugation lemma on the product-form
lemma uses its product formula only; there is no subsequent use of its
cross-basis assertion.\footnote{The corresponding declarations are
\path{tildeObs_selfConsistent_ofGlobalPairWitness},
\path{swapUnitary_conj_tildeObs}, \path{swapUnitary_conj_tildeM},
\path{exists_extractionWitness}, and the two Pauli soundness theorems.
A caller search finds no use of \path{tildeObs_twisted_commutation} outside
its declaration.}

The correction is minimal in the mathematical statement: it restores the
phase determined by the existing observables, preserves the equal-index branch,
and changes no construction, hypothesis, error term, or extraction conclusion.
Retaining the false proposition is a record of the source, not additional
evidence that it is true.  Independent adoption review remains required by
the project protocol.

\section{Conclusion}

The source's cross-basis commutation claim is false for general
$\widetilde u$ and $\widetilde v$.  The correct relation retains the phase
$(-1)^{\operatorname{tr}(e_je_{j'}(\widetilde u\cdot\widetilde v))}$.
Because the later argument does not invoke the false dichotomy, this is a
local correction rather than an additional hypothesis on the extraction
theorem.

\bibliographystyle{alpha}
\bibliography{references}

\end{document}
